\documentclass{article}
\usepackage{amsmath}
\usepackage[hidelinks]{hyperref}
\title{A live observation sent before the split}
\date{}
\begin{document}
\maketitle

\paragraph{Scope and authority.}
Use the R6 planar cut from the pinned
\href{https://arxiv.org/html/2609.29939v1}{KTH managed-removal}
simulation and the exact conditional recovery of the earlier
local-sensing test. Its two cut states differ by $\pm0.8$ of
remote group displacement along a specified lateral direction.
The moving relay that will stop still has a direct view of the
nearest remote robot and a geometric link to the singleton
before the cut. The old node may send a named result while
it is still active. The receiver can choose one of two
shared, preverified bounded control plans, but may not act
as if a required observation had been delivered when absent.
The recovery law is indexed by the known reference cut, two
admitted cut states, named peer, common calibrated frame, caps
and current endpoint. It says nothing about arbitrary earlier
fault times or a different underlying route.
\verb|SOURCE.tsv| pins the physical extraction and the stock
selected SubT dispatch, rolling budget and application callback.
The sensor error, sample/processing delay, link seam, and retained
endpoint scope are authored. The receiver uses a current named
callback, not a ground-truth branch, but this is not an
authentication protocol or a hardware measurement.

\paragraph{A causal certificate with a physical witness.}
A relative-position read in the declared frame measures the
remote position minus the nominal moving reference projected
onto the lateral axis. Use a bounded signed error $\epsilon=0.1$.
The departing node's motion is identical to its nominal
reference in these histories, so its relative position cancels
in this projection. The cutoff correction is at most
$(V+\|v_{\rm nom}\|)\Delta$ for elapsed time $\Delta$ under
speed cap $V=1$. Therefore a decoder can exclude any
candidate cut displacement which lies outside the radius
$\epsilon+(V+\|v_{\rm nom}\|)\Delta$ about the observed value.
This is a conservative, physically justified exclusion, not a
claim that every point in the interval is reachable.

For a late sample ($t=-1$) the nominal source-approved motions
already have opposite offsets of magnitude about $0.5926$.
Even charging the full $0.1$ error and remaining motion,
the opposite $0.8$ alternative at fault time is impossible
in either branch. The relay spends one sample interval,
and the first transmission occurs two quarter ticks after capture.
The controller can choose the checked warm-state motion on a
delivered, correctly scoped code before the fault.
All admissible earlier sender/receiver links and clearance are
checked on the actual reconstructed relative positions with a
continuous chord envelope. Youngness of a received packet alone
is not a premise of the exclusion.

For the early case ($t=-3$), the same positive measured offset
$0.2$ is possible with both future signs. These are not merely
two abstract rows left alive by an overestimate. Both realizable
pre-cut branches drift with the native velocity and add a smooth
cubic translation of the whole remote group from $0.2$ to
$+0.8$ or $-0.8$. The two branches coincide at the sample;
each respects the cap and named old links and clearance.
The first emitted raw-sample sign is still genuinely positive
but one later motion changes sign. Treating that sign as a
cut certificate can direct the left robot into the wrong
recovery, verified by the incompatible realized terminal
positions at the actual bridge distance. Refusing to guess is
not coded as successful mission completion.

\paragraph{Transmission decisions and matched replay.}
The application result contains the branch sign and its modeled
scope in a three-byte word. The earlier erased four-byte debit
can block the first scheduled send while a second timed send
fits once it expires. Conversely, with no earlier debit, a failed
three-byte transmission can itself consume enough TX capacity
that a later physically available path cannot carry the retry.
A wait-for-the-later-slot policy preserves room for that
particular history, but loses another history in which only
the earlier path works. Knowing at least one slot is physically
live without knowing which does not guarantee transport with
these two attempted sends in the same five-byte window.
The first failed send uses three bytes, and the second needs
three; waiting skips the only live slot in the opposite case.
The same truthful evidence and capacity are offered to both
direct and certificate paths. One independent radio callback
succeeds in each fixture.

The finite sweep covers 48 cases per arm per read regime
(selected errors $-0.08,0,0.08$, branch, each live path, and
retained transmit debt). The late-read certificate, equally
informed direct comparator and waiting schedule each meet
24 deadlines; the first two agree trace for trace and the waiting
schedule succeeds on different histories. The early raw control
attempts incompatible motions in 12 cases while a typed check
withholds advance on all unqualified early histories. These
are exact case counts, not estimated scan or delivery frequencies.
The experiment gives an actual modeled way to obtain the
early distinction in part of the declared envelope, and identifies
a separate evidence-expiry law and live capacity obstruction.
It does not implement the original
\href{https://en.endo-lab.jp/research/detail/499}
{Maeda--Endo--Matsuno local crash controller} or establish
a full R6, R3a, independent mission or hardware result.
\end{document}
